% !TeX program = xelatex
% !TeX TXS-program:compile = txs:///xelatex
%

%===============================================================================
% CV – Template (German)
% Content only. Layout/fonts handled by common/cls/cv/*.cls.
%===============================================================================

% documentclass[debug(optional), singlsided/doublesided (2 Page highlightbar left vs. 2 Page highlightbar right ), papersize, global font size]{path to main_cv.cls}
\documentclass[debug,
singlesided,
paper=a4,
fontsize=10pt]{../../common/cls/cv/main_cv}

\usepackage[ngerman]{babel}

%--- options : de/en
\cvSetLanguage{de}

\usepackage{lipsum}

% Skill bar renderer: classic | segmented
\cvUseSkillBarStyle{segmented}


% Hide link borders/colors globally (clean CV look).
% Header links are colored white inside \makeheader via scoped \hypersetup.
\AtBeginDocument{\hypersetup{hidelinks}}

%===============================================================================
% GEOMETRY
%===============================================================================

\setlength\highlightwidth{8cm}
\setlength\headerheight{4cm}
\setlength\marginleft{1cm}
\setlength\marginright{\marginleft}
\setlength\margintop{1cm}
\setlength\marginbottom{1cm}

%===============================================================================
% COLORS
%===============================================================================

% Override shared color defaults (see shared_adjustment.cls for full palette).
%\colorlet{headerbarcolor}{black}
%\definecolor{highlightbarcolor}{HTML}{EAEDED}
\colorlet{headerfontcolor}{white}
\definecolor{highlight}{HTML}{0055FF}
%\colorlet{accent}{highlight}
%\colorlet{heading}{black}
%\colorlet{body}{black}

%===============================================================================
% FOOTER
%===============================================================================

% \makecvfooter{date}{label-page1}{label-page2}{name}
%   Single page: L=date, C=label1, R=name
%   Two pages:   L=date, C=label1/label2, R=Seite n von 2
% \makecvfooter{date}{label-page1}{label-page2}{name}
%   name = {} defaults to \prefix + \name (auto-detected)
\makecvfooter
  {\today}
  {Tabellarischer Lebenslauf}
  {Detaillierte Arbeitserfahrung}
  {}


% --- Personal data ---
\prefix{Dr.}
\name{Max Mustermann}
% --- Contact block ---
% Keep contact rows inside \tagline for a single, reusable header API.
\tagline{
	% Address format: {street and number}{zip code}{city}
	\address{Hauptstraße 101}{D-12345}{Stadt}{Deutschland}
	\phone{+49}{0150 12345678}
	\email{max\_mustermann@email.de}
	\github{sm4rtm4art}{https://github.com/sm4rtm4art}
	\linkedin{Max Mustermann}{https://www.linkedin.com/company/overleaf/?originalSubdomain=uk}
}

%% Photo options => square | roundedbottom
\photo[roundedbottom]{../../images/image}{\dimexpr\headerheight-\marginbottom}


%===============================================================================
% DOCUMENT
%===============================================================================

\begin{document}

	\makeheader%


%===============================================================================
% PAGE 1 — HIGHLIGHT SIDEBAR
%===============================================================================


	\highlightbar{%

		\sidebarSection{Kenntnisse}

		% \cvSideSectionFormat: [Icon]{Content}
		\cvSideSection[\faTerminal]{Programmiersprachen}
		% icons can be found at https://ctan.org/pkg/fontawesome5
		% \cvSkill{skill}{score}
		\cvSkill{\LaTeX}{5}
		\cvSkill{Sprache 1}{4}
		\cvSkill{Sprache 2}{3}
		\cvSkill{Programmiersprache 3}{2}
		\cvSkill{Programmiersprache 4}{1}

		\vspace{0.5cm} %optional

		\cvSideSection[\faChartLine]{Domain 1 }
		%\cvSideSection{Domain 1}

		% \cvSkill{Base}{Detail}
		\cvSkillDetail{Beispiel Skill Section}{5}{Beispiel1· Beispiel2 · Beispiel3}
		\cvSkillDetail{Beispiel Skill Section}{4}{Beispiel1· Beispiel2 · Beispiel3}
		\cvSkillDetail{Beispiel Skill Section}{3}{Beispiel1· Beispiel2 · Beispiel3}
		\cvSkillDetail{Beispiel Skill Section}{2}{Beispiel1· Beispiel2 · Beispiel3}

		\vspace{0.5cm} %optional

		\cvSideSection[\faDatabase]{Domain 2}
		%\cvSideSection{Domain 2}
		\cvSkillDetail{Beispiel Skill Section}{5}{Beispiel1· Beispiel2 · Beispiel3}
		\cvSkillDetail{Beispiel Skill Section}{4}{Beispiel1· Beispiel2 · Beispiel3}
		\cvSkillDetail{Beispiel Skill Section}{3}{Beispiel1· Beispiel2 · Beispiel3}
		\cvSkillDetail{Beispiel Skill Section}{2}{Beispiel1· Beispiel2 · Beispiel3}

		\vspace{0.5cm} %optional

		\cvSideSection[\faGlobe]{Sprachen}
		%\cvSideSection{Sprachen}
		\cvSkill{Deutsch / Englisch}{5}
		\cvSkill{Französisch}{3}

		\vspace{0.5cm} %optional

		\sidebarSection{Zertifikate}
		% \cert{Certificate}{Detail}
		\cert{Zertifikat 1}{Zertifikatdetail}
		\cert{Zertifikat 2}{Zertifikatdetail}
		\cert{Zertifikat 3}{Zertifikatdetail}
		\cert{Zertifikat 4}{Zertifikatdetail}
		\cert{Zertifikat 5}{Zertifikatdetail}
		\cert{Zertifikat 6}{Zertifikatdetail}


	}%


	%===============================================================================
	% PAGE 1 — MAIN CONTENT
	%===============================================================================

	\mainbar{

		\section[\faUser]{Über mich}
		%\section{Über mich}

		 \lipsum[1][1-8]

		\vspace{1.0em}

		\section[\faBriefcase]{Arbeitserfahrung}
		%\section{Arbeitserfahrung}

		% Unified Format: {Company}{Position}{Date}{Location}{Highlighttext}
		\cvJobOverview{Firma}
		{Position 1 | Position 2}
		{12|25 \textendash Aktuell} % ATS- Info => 12/2025 - \today{mm/yyyy}
		{Ort 1}
		% No highlight — optional [highlight] argument omitted

		\cvJobOverview{Firma 2}
		{Position 1 }
		{01|25 \textendash 11|25}  % ATS- Info => 01/2025 - 11/2025
		{Ort 2}
		[{\textbf{Optional:} Highlight Text, der die Arbeitsbeschreibung definiert.}]

		\cvJobOverview{Firma 3}
		{Position 1 }
		{02|23 \textendash 12|24}
		{Ort 3}
		[{\lipsum[1][1-3]}]

		\vspace{1cm} % Optional

		\section[\faGraduationCap]{Ausbildung}
		%\section{Ausbildung}

		% Unified Format: {Institution}{Degree}{Date}{Location}{Details}

		\cvEduOverview{Ausbildungsstätte 1}
		{Abschluss}
		{01|2022 \textendash 02|2023}  % ATS- Info => 10/2022 - 02/2023
		{Ort}
		% No highlight — optional [highlight] argument omitted

		\cvEduOverview{Ausbildungsstätte 2}
		{Abschluss}
		{01|2020 \textendash 12|2021}
		{Ort}
		[{\textbf{Optional:} Note, Abschluss, Details}]

		\cvEduOverview{Ausbildungsstätte 3}
		{Abschluss}
		{01|2017 \textendash 12|2019}
		{ Längerer Ortsname }
		[{\textbf{Optional:} Note, Abschluss, Details}]

	}%

	\makebody%
	\clearpage   % Important for creating the new page

	%===============================================================================
	% PAGE 2
	%===============================================================================

	\pagestyle{highlightmain} % Sets style

	%===============================================================================
	% PAGE 2 — HIGHLIGHT SIDEBAR
	%===============================================================================


	\highlightbar{%
		\sidebarSection{Methodenprofil}


		\cvSideSection{Methodik 1}
		\begin{cvTagCloud}
			\cvSkillTags{
				Skill 1,
			Skill 2,
			Skill 3,
			Skill 4,
			Skill mit längerer Zeichenlänge
			}
		\end{cvTagCloud}

		\vspace{0.2cm}

		\cvSideSection{Methodik 2}
			\begin{cvTagCloud}
				\cvSkillTags{
					Skill 1,
					Skill 2,
					Skill 3,
					Skill 4,
					Skill mit längerer Zeichenlänge
				}
			\end{cvTagCloud}

		\vspace{0.2cm}

		\cvSideSection{Methodik 3}
			\begin{cvTagCloud}
				\cvSkillTags{
					Skill 1,
					Skill 2,
					Skill 3,
					Skill 4,
					Skill mit längerer Zeichenlänge
				}
			\end{cvTagCloud}
	}%



  	%===============================================================================
  	% PAGE 2 — MAIN CONTENT
  	%===============================================================================

	\mainbar{

		%\section{Arbeitserfahrung im Detail}
		\section[\faBriefcase]{Arbeitserfahrung im Detail}

		%---- \cvJobDetails

		% 1. Company 1 –{Company}{Position}{Date}{Location}{Job description}{Tags}
		\cvJobDetail%
		{Firma 1}
		{Position 1.1 | Position 1.2 }
		{12|2025\textendash Aktuell}
		{Ort1}
		{
		% ---  optional  {Core} / short description ---
		\cvDetailItem{Optional:} {Details aus der Arbeitserfahrung die man hier beschreiben kann}
		\cvDetailItem{Optional:} {\lipsum[1][1-3]}
		\cvDetailItem{Optional:} {\lipsum[1][1-4]}
		}
		% No tags — optional [tags] argument omitted

		% 2. Company 2 –{Company}{Position}{Date}{Location}{Job description}{Tags}

		\cvJobDetail%
		{Firma 2}
		{Position 2 }
		{01|2025 \textendash 11|2026}
		{Ort 2}
		{
			% ---  optional  {Core} / short description --
			\cvDetailItem{Optional:} {Details aus der Arbeitserfahrung die man hier beschreiben kann}
			\cvDetailItem{Optional:} {Details aus der Arbeitserfahrung die man hier beschreiben kann}
		}
		[{% --- Tags (pick a shorter set if needed) ---
			\cvJobTags{
				Tag 1,
				Tag 2,
				Tag 3,
				Tag 4  as pill form,
				Tag 5  as pill form,
				Tag 6  as pill form
				}
		}]

		% 3. Company 3 –{Company}{Position}{Date}{Location}{Job description}{Tags}
		\cvJobDetail%
		{Firma 3}
		{Position 3}
		{02|2023\textendash12|2024}
		{Ort 3}
		{
		% ---  optional  {Core} / short description ---
		\cvDetailItem{Optional:} {Details aus der Arbeitserfahrung die man hier beschreiben kann}
		\cvDetailItem{Optional:} {Details aus der Arbeitserfahrung die man hier beschreiben kann}
		}
		[{% --- Tags (pick a shorter set if needed)
			\cvJobTags{
			Tag 1,
			Tag 2,
			Tag 3,
			Tag 4  as pill form,
			Tag 5  as pill form,
			Tag 6  as pill form
			}
		}]


}%



\makebody%

\end{document}
